\songtitle{Morocho}

\begin{notes}
Based on a characters from \emph{Ricochets} by Carlos Thompson.\\
Character: Miguel Morocho.\\
Cowritten with ChatGPT based on the arc in the novel, the companion texts, and further input by the author.
\end{notes}

\begin{style}
salsa, piano tumbao, sharp brass hits, güiro percussion, montuno call and response, male spoken intro, dramatic tenor lead, choral responses, minor key, 145 BPM, live conga groove, timbales breaks, walking bass lines, trumpet stabs, crowded mix, plate reverb, punchy snare, tension release, grim revenge, urgent momentum
\end{style}

\begin{lyrics}

\songpart{Intro}
\annotation{piano enters with a mysterious tumbao, the Caribbean güiro and the brass instruments with a sharp hit}

\annotation{spoken}\annotation{piano only}\\
Miguel Morocho era un tipo astuto\\
tal vez demasiado astuto y esa fue su perdición

\songpart{Verse 1}
Miguel Morocho un buen día, en inteligencia entrenó,\\
gringa instrucción recibió siendo de la policía.\\
Mas luego, por otra vía,\\
al de Pereira cartel sirvió de asesor aquel.\\
Dos abogados en Miami le daban un gran tsunami,\\
y con Carlson bebió hiel.

\songpart{Pre-Chorus}
Con su instructor se sentó,\\
Sarah el blanco resultaba;\\
entre copas se pactó\\
lo que el cartel tanto ansiaba.

\songpart{Chorus}
Morocho era muy astuto, de policía al cartel,\\
tejiendo redes de sangre con hilos de fría hiel.\\
Mas no hay pacto que perdone ni astucia que venza al mal:\\
la sangre de aquella niña lo llevó a su funeral.

\annotation{Instrumental Break, güiro and brass}

\songpart{Verse 2}
Guerra civil de carteles, sangre y fuego en la región;\\
Mojica busca una unión vistiendo falsos laureles.\\
A Alcántara da timoneles mientras se desmoviliza,\\
mas la paz se minimiza: limpiando casa los dos,\\
Morocho al de la DEA dio la red que ya no utiliza.

\songpart{Chorus}
Morocho era muy astuto, de policía al cartel,\\
tejiendo redes de sangre con hilos de fría hiel.\\
Mas no hay pacto que perdone ni astucia que venza al mal:\\
la sangre de aquella niña lo llevó a su funeral.

\songpart{Montuno}
\annotation{Call and Response, energetic, call and response style}\\
\annotation{Choir: multiple voices} (¿Y quien era aquella niña?)\\
\annotation{Solo: lead voice} Era la hija de Carlson\\
\annotation{Choir:} (¿Y qué le pasó a la niña?)\\
\annotation{Solo:} Con su mamá la mataron

\annotation{Choir:} (¿Y qué le pasó a don Carlson?)\\
\annotation{Solo:} A Bogotá fue a parar.\\
\annotation{Choir:} (Eso no lo entiendo, hermano)\\
\annotation{Solo:} Que su pena se fue a expiar.

\songpart{Bridge}
Morocho lo halla en las sombras,\\
buscando su redención;\\
ya Carlson rechaza el trago\\
curando su perdición.

Le ofrece armas y recursos\\
para saciar su venganza,\\
si a Bard le lleva un correo\\
como pacto de confianza.

\songpart{Verse 3}
Bard acepta aquel convenio y limpia casa el cartel,\\
mientras se cobra la piel de socios en el milenio.\\
Mas Carlson, con gran ingenio, ve al enemigo real;\\
cayó la red de metal por el abogado vivo,\\
y con soplón expeditivo la DEA borró el dineral.

\songpart{Chorus}
Morocho era muy astuto, de policía al cartel,\\
tejiendo redes de sangre con hilos de fría hiel.\\
Mas no hay pacto que perdone ni astucia que venza al mal:\\
la sangre de aquella niña lo llevó a su funeral.

\songpart{Outro}
Carlson fue un fantasma,\\
un destello que se deshizo.\\
Luego, el cuerpo:\\
Morocho en su casa,\\
la femoral cortada,\\
herida íntima y cercana,\\
y su arma inútil a un paso.

Ernesto lee la escena sin pánico;\\
vacia memorias, legajos, secretos,\\
reparte el oro y las joya\\
se queda el poder.


\end{lyrics}

